\begin{exercise}[Wilson loops as local shifts of the theta angle]{I-CH10-011}
In the two-dimensional charge-$m$ Higgs model, suppose a unit instanton has
flux $\int F=2\pi/m$.  Write its dilute-gas fugacity per unit area as $\zeta$,
so that on an area $\mathcal V$
\[
  Z(\mathcal V,\theta)=
  \exp\left[2\zeta\mathcal V\cos\left(\frac\theta m\right)\right].
\]
For a charge-$k$ Wilson loop $W_k(C)=\exp(\ii k\oint_C A)$ enclosing area
$A$, use Stokes' theorem to find the phase contributed by every instanton or
anti-instanton inside $C$.  Prove that the interior gas is obtained by
$\theta\mapsto\theta+2\pi k$.  Split the gas into the inside and outside
regions, divide by the vacuum partition function, and derive
\[
  \langle W_k(C)\rangle=
  \exp\left\{2\zeta A\left[
  \cos\left(\frac{\theta+2\pi k}{m}\right)
  -\cos\left(\frac\theta m\right)\right]\right\}
\]
at dilute-gas order.  State when its area coefficient vanishes.
\end{exercise}
